\documentclass[12pt, titlepage]{article}
%\usepackage[norsk]{babel}
\usepackage[utf8]{inputenc}
\usepackage{minted, enumerate}
\usepackage[colorlinks = true,
linkcolor = blue,
urlcolor = blue]{hyperref}


\begin{document}

\title{Mappeoppgave - del 2}
\author{IIRA2001 - Høst 2023}
\date{}
\maketitle


Andre del av mappen dreier seg om dataanalyse og web-APIer.

Du skal legge ved 2 programmer hvor du gjør noe interessant med 1 eller flere datasett.
Programmene kan undersøke en problemstilling, utforske data, teste en hypotese eller gjøre noe annet nyttig.
Vi vil gjerne at du i ett av programmene bruker et web-api eller et pythonbibliotek (eksempelvis yfinance) til å hente dataene du trenger fra nettet.
Eventuelt kan 1 av programmene basere seg kun på bruk av et web-api og ikke dataanalyse.


\section*{Oppgave}
Skriv to pythonprogrammer hvor du:
\begin{enumerate}[$\bullet$]
   \item {analyserer 1 eller flere datasett}
   \item {gjerne i et 1 av programmene bruker et web-api eller bibliotek til å hente data programmatisk}
\end{enumerate}

\subsection*{Unntak til punktene over}
\begin{enumerate}[$\bullet$]
   \item {Dersom du har en kul idé til et program som bruker et web-api, men som ikke faller innen kategorien dataanalyse kan dette sikkert godkjennes, men hør med foreleser for sikkerhets skyld}
   \item {Dersom du har to gode dataanalyse-programmer som ikke henter inn data programmatisk fra et web-api (eller yfinance), godkjennes dette mot et \textit{mindre} trekk under vurderingen}
\end{enumerate}

\section*{Oppgavetips og eksempler}

Du kan, litt som i første del av mappen, ta utgangspunkt i øvingsoppgaver og eksempler fra timene.

\subsection*{Utforske data}

I timen lastet vi inn et datasett fra SSB med populasjonsdata fra kommunene i Norge og plottet endring i populasjonsstørrelse for noen kommuner.

Fraflytting fra Nord-Norge har lenge vært en bekymring (se feks. \href{https://www.nrk.no/tromsogfinnmark/4000-flyttet-fra-nord-norge-i-2022-1.16345861}{denne} artikkelen fra nrk),
dette kunne vært et utgangspunkt for et pythonprogram som utforsker populasjonsdata. Du kan for eksempel plotte grafer for netto innflytning, populasjonsstørrelse i ulike regioner og se om man finner denne fraflytningstrenden.
Kanskje kan det være interessant å sette dette i sammenhengen med aldersfordelingen også.

En god oppgave kommenterer selvfølgelig også funnene.

\subsection*{Problemstilling}

I eksempelet over har vi kanskje ikke noe spørsmål vi vil finne svar på. I forskning kalles et spørsmål man vil finne svar på en problemstilling. 
Vi har i timen hentet data om arbeidsledighet og åpnede konkurser, og undersøkt om det er en sammenheng mellom de to. Vi har også undersøkt om det er en sammenheng
mellom fraflytting fra en kommune og arbeidsledigheten/sysselsettingen i en øving.
Man kan se etter (lineære) sammenhenger mellom to statistiske variabler ved å regne ut korrelasjonen mellom de, slik som i timen. Man kan også plotte grafer, prikkeplott eller histogram
å se etter sammenhenger.

Det finnes mange interessante ting å undersøke om henger sammen, feks: Strømpris og prisen på aluminium, strømpris i Europa og strømpris i Norge, styringsrente og inflasjon,
byggevareindeks og huspriser osv.

\subsection*{Hypotesetesting}

Fra statistikken har dere lært å teste hypoteser. Mulige hypoteser å med data fra SSB teste er "folk med høyere utdanning har høyere lønn". Eksempel er planlagt å gi i en øving eller forelesning

\subsection*{Nyttige program}

Du kan også lage et program av praktisk nytte, eller som du synes er artig. Vi har eller vil gjennomgå et program som regner om mellom ulike valutaer med ferskest mulige kurser.
Den ene øvingsoppgaven dreier seg om å hente quiz-spørsmål fra opentdb.com sitt web-api og kjøre en liten quiz. 


\section*{Øvrig innhold}

I tillegg til pythonprogrammene, skal det leveres en slags rapport.
Vi vil ha beskrevet hva målet med hvert program er, hva ideen var, eller hva problemstillingen er.

I tillegg trenger vi en del hvor dere reflekterer og diskuterer resultatene. Hva har du funnet ut? Hvordan fungerer programmet ditt?

Dette er ikke slik som refleksjonsnotatet fra mappe del 1, men mer en rapport angående hvert program.

Hvert av programmene trenger en egen rapport



\section*{Tilbakemelding}

Siste undervisningsuke vil det bli arrangert medstudentvurdering i grupper. Hver gruppe velger ut en av besvarelsene til å få en tilbakemelding alle i gruppen kan se.
Det beste hadde naturligvis vært å gi alle individuell tilbakemelding, men det vil vi ikke kunne rekke før innlevering

\section*{Innlevering}

Innlevering skjer i Inspera. Det blir åpnet for innlevering 18. desember, og stenges 22. desember. Hele mappen må altså være ferdig og leveres mellom 18-22. desember.
Mer informasjon kommer

\end{document}
